\section{Research Questions}
\label{sec:rq}

Scientific output grows faster than any researcher can read. Turning that
literature into \emph{what has not yet been done}---concrete, novel, feasible
research directions---is a core scholarly task that does not scale manually, and
that na\"ive per-paper LLM use cannot scale \emph{economically}: a per-paper
extractor costs on the order of \$4{,}000 per million papers and is rate-limited.
This thesis builds and evaluates \textbf{Gap2Idea}, an end-to-end system that
mines research gaps from papers, organises them into a graph, and generates
evaluated research ideas. Three research questions structure both the system and
the literature; the pipeline instantiates them in sequence---\textbf{RQ2
(extract) $\rightarrow$ RQ1 (graph) $\rightarrow$ RQ3 (ideate)}.

\subsection{RQ2 --- Text-Based Gap Mining}
\emph{Can research gaps (limitations, future-work directions) be extracted from
scientific papers at corpus scale without a per-paper LLM, at a quality
competitive with fine-tuned transformers?}
\begin{itemize}
  \item \textbf{RQ2a.} What architecture localises and classifies gap sentences
        cheaply and at high recall? (A structural slice $\rightarrow$ a cheap
        classifier $\rightarrow$ an LLM precision filter.)
  \item \textbf{RQ2b.} Is extraction quality bounded by \emph{training data} or
        \emph{model capacity}?
  \item \textbf{RQ2c.} How does the funnel compare, on public benchmarks (LimGen,
        BAGELS, the future-work corpus), to fine-tuned and generative baselines?
\end{itemize}

\subsection{RQ1 --- Graph-Based Forecasting}
\emph{Can research gaps be organised into a multi-relational graph whose structure
identifies promising, novel directions---by recombination and bridging---better
than ungrounded pairing?}
\begin{itemize}
  \item \textbf{RQ1a.} How to embed and cluster gaps into interpretable
        communities?
  \item \textbf{RQ1b.} Do graph-derived \emph{bridge} and \emph{frontier} seeds
        yield better idea candidates than random gap pairs? (The extrinsic
        clustering signal.)
\end{itemize}

\subsection{RQ3 --- LLM-Driven Ideation}
\emph{Can an LLM, grounded in retrieved gap evidence, generate novel, feasible,
non-hallucinated research ideas, and can idea quality be evaluated reliably?}
\begin{itemize}
  \item \textbf{RQ3a.} What grounding and anti-hallucination constraints keep
        generated ideas faithful to evidence?
  \item \textbf{RQ3b.} Does a multi-agent critic/revise loop improve idea quality
        over a single pass?
  \item \textbf{RQ3c.} Can idea quality be assessed by combining an LLM-judge
        panel with a human expert study, and do LLM judgments correlate with human
        ratings?
\end{itemize}
